% Part: turing-machines
% Chapter: undecidability
% Section: universal-tm

\documentclass[../../../include/open-logic-section]{subfiles}

\begin{document}

\olfileid{tur}{und}{uni}
\olsection{সার্বজনীন টুরিং যন্ত্র}

\olref[enu]{sec}-এ আমরা আলোচনা করেছি, কীভাবে প্রতিটি টুরিং যন্ত্রকে
পূর্ণসংখ্যার একটি সসীম অনুক্রম দিয়ে বর্ণনা করা যায়। এই অনুক্রমটি টুরিং
যন্ত্রের দশা, বর্ণমালা, আরম্ভিক দশা ও নির্দেশগুলির সংকেতায়ন করে। এমন সব
বর্ণনার সেটটি যে !!{enumerable}, সে কথাও আমরা উল্লেখ করেছি। এমন বর্ণনার
সেটটি !!{denumerable} বলে~$\Nat$ থেকে এই বর্ণনাগুলিতে !!a{surjective}
অপেক্ষক আছে। যেমন, কান্টরের আঁকাবাঁকা পদ্ধতি ব্যবহার করে এমন
!!a{surjective} অপেক্ষক পাওয়া যায়। এটি সব টুরিং যন্ত্রের (বর্ণনা)
তালিকায়িত করার একটি উপায় দেয়। এমন একটি তালিকায়ন স্থির করলে এবার
$1$-ম, $2$-য়, \dots, $e$-তম টুরিং যন্ত্রের কথা বলার অর্থ হয়; শেষটির
সূচক~$e$। এই সংখ্যাগুলিকে \emph{সূচক} বলা হয়।

\begin{defn}
$M$ যদি (আমাদের স্থির তালিকায়নে) $e$-তম টুরিং যন্ত্র হয়, তবে~$e$-কে
~$M$-এর একটি \emph{সূচক} বলি। $e$-তম টুরিং যন্ত্রকে $M_e$ লিখি।
\end{defn}

একটি যন্ত্রের একাধিক সূচক থাকতে পারে। যেমন,~$M$-এর দুটি বর্ণনায় তার
নির্দেশগুলি যে ক্রমে তালিকায় লেখা হয় শুধু সেটিই আলাদা হতে পারে, আর এই
ভিন্ন বর্ণনাদুটির সূচকও ভিন্ন হবে।

গুরুত্বপূর্ণ বিষয় হলো, টুরিং যন্ত্রের বর্ণনার তালিকায়নটি এমনভাবে দেওয়া
সম্ভব যাতে সূচক থেকে~$M$-এর বর্ণনা কার্যকরভাবে গণনা করা যায় এবং কোনো
যন্ত্র~$M$-এর বর্ণনা থেকে তার একটি সূচক কার্যকরভাবে গণনা করা যায়। তখন
চার্চ--টুরিং থিসিস অনুসারে এমন একটি টুরিং যন্ত্র পাওয়া সম্ভব, যা সূচক
~$e$-এর টুরিং যন্ত্রটির বর্ণনা উদ্ধার করে এবং সংশ্লিষ্ট বর্ণনাটি নির্গম
হিসেবে নিজের ফিতায় লেখে। বর্ণনাটি হবে~$\TMstroke$-এর কয়েকটি খণ্ডের
অনুক্রম (যেগুলি~$M_e$-র বর্ণনাকারী অনুক্রমের ধনাত্মক পূর্ণসংখ্যাগুলিকে
প্রকাশ করে)।

এই ফল থাকলে স্বাভাবিকভাবেই প্রশ্ন ওঠে: টুরিং যন্ত্রের সূচকের কোন
অপেক্ষকগুলি নিজেরাই টুরিং যন্ত্রে গণনসাধ্য? টুরিং যন্ত্রের সূচকের কোন
বৈশিষ্ট্যগুলি টুরিং যন্ত্রে নির্ণয় করা যায়? একটি উদাহরণ: সূচক~$e$-কে ওই
সূচকের টুরিং যন্ত্রের দশার সংখ্যায় পাঠানো অপেক্ষকটি টুরিং যন্ত্রে
গণনসাধ্য। এমন একটি টুরিং যন্ত্র যা করবে তা হলো: ফিতায় $e$টি
~$\TMstroke$-এর একটিমাত্র খণ্ড নিয়ে শুরু করে, প্রথমে $e$-কে তার বর্ণনায়
বিসংকেতায়িত করবে। এবার ফিতায়~$\TMstroke$-এর কয়েকটি খণ্ডের অনুক্রম দিয়ে
বর্ণনাটি প্রকাশিত হয়। এই অনুক্রমের প্রথম !!{element} যেহেতু দশার সংখ্যা,
তাই এখন প্রথম~$\TMstroke$-খণ্ডটি ছাড়া বাকি সব মুছে থামলেই হবে।

নিচের ফলটি বিশেষ উল্লেখযোগ্য:

\begin{thm}\ollabel{thm:universal-tm} এমন একটি \emph{সার্বজনীন টুরিং
  যন্ত্র}~$U$ আছে, যা নিবেশ $\tuple{e,n}$-এ শুরু করলে
  \begin{enumerate}
    \item $M_e$ নিবেশ~$n$-এ থামলে এবং কেবল তখনই থামে, এবং
    \item $M_e$ নির্গম $m$ নিয়ে থামলে~$U$-ও তাই করে।
  \end{enumerate}
  অতএব~$U$ অপেক্ষক $f\colon \Nat \times \Nat \pto \Nat$ গণনা করে, যেখানে
  $M_e$ নিবেশ~$n$-এ শুরু করে নির্গম~$m$ নিয়ে থামলে $f(e,n) = m$, এবং
  অন্যথায় অপেক্ষকটি অসংজ্ঞায়িত।
\end{thm}

\begin{proof}
  যন্ত্র~$U$ বাস্তবে তৈরি করা মূলত অসম্ভব, কারণ যন্ত্রটি অত্যন্ত জটিল।
  কিন্তু সেটি কীভাবে কাজ করে তার একটি রূপরেখা বর্ণনা করে পরে চার্চ--টুরিং
  থিসিস আহ্বান করতে পারি। শুরুতে~$U$-এর ফিতায় $e$টি~$\TMstroke$-এর একটি
  খণ্ডের পরে $n$টি~$\TMstroke$-এর একটি খণ্ড থাকে। সেটি প্রথমে নিবেশ~$n$-এর
  ডান দিকে সূচক~$e$-কে “বিসংকেতায়িত” করে। এতে~$M_e$-এর নির্দেশ বর্ণনাকারী
  সংখ্যার একটি তালিকা (অর্থাৎ~$\TMblank$ দিয়ে আলাদা করা~$\TMstroke$-
  খণ্ডগুলি) পাওয়া যায়। এরপর~$U$ ওই~$M_e$-র বর্ণনার ডান দিকে~$M_e$-র
  আরম্ভিক দশার সংখ্যা এবং সংখ্যা~$1$ লেখে। (আবার, এগুলি এককীয় উপস্থাপনায়
  ~ $\TMstroke$-এর খণ্ড দিয়ে প্রকাশিত।) তারপর সেটি নিবেশটি ($n$টি
  ~ $\TMstroke$-এর খণ্ড) ডান দিকে অনুলিপি করে—তবে প্রতিটি~$\TMstroke$-এর
  বদলে তিনটি~$\TMstroke$-এর একটি খণ্ড বসায় (মনে রেখো,~$\TMstroke$
  প্রতীকের সংখ্যা~$3$; $\TMendtape$-এর সংখ্যা $1$ এবং~$\TMblank$-এর সংখ্যা
  $2$)। (নিজের ফিতায়~$\TMblank$ প্রতীক দিয়ে আলাদা করা) এই খণ্ড-অনুক্রমের
  বাঁ প্রান্তে সেটি একটিমাত্র~$\TMstroke$, অর্থাৎ~$\TMendtape$-এর সংকেতটি
  লেখে।

  এবার~$U$-এর ফিতায় আছে: সূচক~$e$, সংখ্যা~$n$, আরম্ভিক দশার সংকেত-সংখ্যা
  (“বর্তমান দশা”), আরম্ভিক হেড-অবস্থান~$1$-এর সংখ্যা (“বর্তমান হেডের
  অবস্থান”), এবং “ফিতা”-র আরম্ভিক বিষয়বস্তু (ফাঁকা~$\TMblank$ দিয়ে আলাদা
  করা~$M_e$-র প্রতীকগুলির সংকেত-সংখ্যা—“প্রতীকগুলি”—প্রকাশকারী
  ~ $\TMstroke$-খণ্ডের একটি অনুক্রম)।

  এবার নিচের কাজগুলি করে সেটি নিবেশ~$n$-এ শুরু করা~$M_e$-র কাজ অনুকরণ
  করে:
  \begin{enumerate}
    \item “বর্তমান হেডের অবস্থান”-এর সংখ্যা~$k$ খুঁজে বের করো (শুরুতে
    সেটি~$1$),
    \item সেখানে কোন “প্রতীক” আছে তা দেখতে “ফিতা”-র $k$-তম খণ্ডে যাও,
    \item\ollabel{find-inst}%
    বর্তমান “দশা” ও “প্রতীক”-এর সঙ্গে মেলে এমন নির্দেশটি খুঁজে বের করো,
    \item “ফিতা”-র $k$-তম খণ্ডে ফিরে সেখানে থাকা “প্রতীক”-এর জায়গায়
    ~ $M_e$ যে প্রতীকটি লিখত তার সংকেত-সংখ্যা বসাও,
    \item যে জায়গায় বর্তমান “দশা” নথিভুক্ত আছে সেখানে গিয়ে সংখ্যাটির
    জায়গায় নতুন দশার সংখ্যা বসাও,
    \item যে জায়গায় “ফিতার অবস্থান” নথিভুক্ত আছে সেখানে গিয়ে একটি
    ~ $\TMstroke$ মুছে দাও অথবা একটি~$\TMstroke$ যোগ করো (নির্দেশটিতে
    যথাক্রমে বাঁ দিকে অথবা ডান দিকে চলতে বলা থাকলে)।
    \item পুনরাবৃত্তি করো।\footnote{এখানে আমরা কয়েকটি সূক্ষ্ম অসুবিধা
    এড়িয়ে যাচ্ছি। যেমন, “বর্তমান হেডের অবস্থান” যে গণকে~$U$ অনুসরণ করে,
    সেটি বাড়ানোর সময় কিছু অতিরিক্ত জায়গা লাগতে পারে—সেক্ষেত্রে পুরো
    “ফিতা”-টিকে ডান দিকে সরাতে হবে।}
  \end{enumerate}
  নিবেশ~$n$-এ শুরু করা~$M_e$ কখনও না থামলে~$U$-ও কখনও থামে না, ফলে তার
  নির্গম অসংজ্ঞায়িত।

  ধাপ~\olref{find-inst}-এ যদি দেখা যায় যে~$M_e$-র বর্ণনায় বর্তমান
  “দশা”/“প্রতীক” যুগলের কোনো নির্দেশ নেই, তবে~$M_e$ থামত। এমন হলে~$U$
  নিজের ফিতার “ফিতা”-অংশের বাঁ দিকের অংশটি মুছে দেয়। এরপর সংকেতায়িত
  “ফিতা”-র প্রথম এক-দাগের খণ্ডটি শেষ-চিহ্নের সংকেত কি না যাচাই করে
  সেটি অতিক্রম করে। তার পরের প্রতিটি তিন~$\TMstroke$-এর খণ্ডের জন্য
  (যা~$M_e$-র ফিতায় একটি~$\TMstroke$ প্রকাশ করে) নিজের ফিতার বাঁ প্রান্তে
  একটি~$\TMstroke$ লেখে এবং ক্রমে সংকেতায়িত “ফিতা” মুছে দেয়। এরপর কোনো
  দুই-দাগের ফাঁকা-প্রতীকের সংকেত থাকলে তার পরের বাকি সব খণ্ডও একই
  ফাঁকা-প্রতীকের সংকেত কি না~$U$ যাচাই করে সেগুলি মুছে দেয়। কাজ শেষ
  হলে~$U$-এর ফিতায় দৈর্ঘ্য~$m$-এর
  ~ $\TMstroke$-এর একটিমাত্র খণ্ড থাকে।

  প্রারম্ভিক এক-দাগের শেষ-চিহ্নের সংকেত না পেলে, নির্গমের অংশে তিন
  ~ $\TMstroke$-এর খণ্ড ছাড়া অন্য কিছু পেলে, অথবা প্রথম দুই-দাগের
  ফাঁকা-প্রতীকের সংকেতের পরে অন্য কোনো সংকেত পেলে~$U$ অবিলম্বে এমনভাবে
  থামে যাতে তার ফিতায়~$\TMstroke$-এর একটিমাত্র খণ্ড না থাকে। তাই এই
  ক্ষেত্রে~$U$-এর নির্গম কোনো স্বাভাবিক সংখ্যা নয়, অর্থাৎ~$f(e,n)$
  অসংজ্ঞায়িত।
% BN-SRC-176 TARGET-CORRECTION-BEGIN
\par\smallskip\noindent\textnormal{\small\textbf{উৎস-সংশোধন (BN-SRC-176):} হিমায়িত ইংরেজি নির্গম-পর্বে সংকেতায়িত “ফিতা”-র প্রতিটি তিন-দাগের খণ্ড অনুবাদ করতে এবং অন্য যেকোনো খণ্ড পেলেই ব্যর্থ হতে বলে। কিন্তু আগেই ওই “ফিতা”-র প্রথম খণ্ডে এক-দাগের শেষ-চিহ্ন রাখা হয়েছে, আর উপস্থিত ফাঁকা প্রতীকের সংকেত দুই দাগের; ফলে বর্ণিত পরীক্ষা প্রতিটি সুগঠিত নির্গমকেই তার প্রথম খণ্ডে প্রত্যাখ্যান করত। বাংলা পাঠে প্রথম শেষ-চিহ্ন যাচাই, মধ্যবর্তী তিন-দাগের নির্গম, এবং থাকলে কেবল পশ্চাৎবর্তী দুই-দাগের ফাঁকা অংশ—এই সম্পূর্ণ সুগঠিত রূপটি যাচাই করা হয়েছে।} % BN-SRC-176 TARGET-CORRECTION-END
\end{proof}

\end{document}
